\section{Kimi K2与Agent能力}

基于Muon优化器和稳定训练技术，Kimi K2在强化学习后训练后取得了多项突破：

\begin{itemize}
\item \textbf{Agent能力全面提升}：可对标美国前沿公司
\item \textbf{HLE（Humanities Last Exam）}：达到45\%准确率，超过OpenAI，这是一个人类也难以解答的高难度Benchmark
\item \textbf{K2 Thinking}：完成连续200--300步的工具调用，中间持续思考、搜索、编写Python程序
\item \textbf{Kimi K2是中国第一个Agentic模型}
\end{itemize}

\begin{knowledgebox}{为什么Agent本质上是搜索问题}
Agent的推理或RL训练本质上是一个搜索过程。例如从头开发一个Linux操作系统——如果有无限compute，可以枚举所有token组合找到正确答案。更好的基础模型提供了更强的先验（heuristic），减少了搜索空间。Token Efficiency的提升等价于更强的先验，long context则提供了更强的工作记忆（环境感知能力）。
\end{knowledgebox}


